% cSpell:locale en,pt-BR - Adicionado para suporte ao corretor ortográfico
% Extensão Code Spell Checker no VS Code (muito útil!)
\chapter{Ilustrações, tabelas e equações}\label{cap:ilustracoes}

\section{Figuras e Tabelas}\label{sec:figuras-tabelas}

Assim como as citações, figuras e tabelas são elementos onipresentes em trabalhos acadêmicos e
técnicos. Eles ajudam a ilustrar conceitos, apresentar dados e facilitar a compreensão do conteúdo.
Figuras e tabelas devem ser inseridas no texto próximo ao local onde são mencionadas pela primeira
vez, para facilitar a leitura e evitar que o leitor precise procurar por elas em outras partes do
documento. Além disso, é importante numerar e rotular corretamente cada figura e tabela, utilizando
os ambientes \texttt{figuraabnt} para figuras e \texttt{tabelaabnt} para tabelas, fornecidos por este
modelo. Isso garante que as referências a esses elementos sejam atualizadas automaticamente, mesmo
que a ordem ou o número de figuras e tabelas mude durante o processo de edição.

A norma ABNT NBR 14724:2024~\cite{nbr14724:2024} estabelece diretrizes específicas para a
formatação e apresentação de figuras e tabelas em trabalhos acadêmicos. De acordo com essa norma,
as figuras devem ser numeradas consecutivamente em algarismos arábicos (Figura 1, Figura 2, etc.) e
devem possuir uma legenda descritiva, localizada acima da figura. As tabelas também devem ser
numeradas consecutivamente em algarismos arábicos (Tabela 1, Tabela 2, etc.) e devem possuir uma
legenda descritiva, localizada acima da tabela. Além disso, é importante garantir que as figuras e
tabelas estejam devidamente referenciadas no texto. Este modelo oferece os comandos
\verb|\reffig{}| e \verb|\reftable{}| para autoreferências formatadas, bastando passar o rótulo da figura/tabela. Use os prefixos \texttt{fig:} para figuras (por exemplo, \texttt{fig:exemplo})
e \texttt{tab:} para tabelas (por exemplo, \texttt{tab:exemplo}) ao criar os rótulos. Também existem
os comandos \verb|\reffigcomp{}| e \verb|\reftablecomp{}| para citar uma figura junto com seu título. Esta não é uma prática padrão, mas pode ser interessante em alguns contextos.

Figuras e tabelas devem referenciar suas fontes também, utilizando o comando \verb|\fonte{}|
fornecido por este modelo, que posiciona a fonte de forma adequada conforme as normas ABNT.  Note
que a fonte deve ser apresentada em fonte menor que o corpo do texto e alinhada à esquerda, conforme
exemplificado a seguir. Caso a Figura ou Tabela sejam de autoria própria, utilize a expressão ``o
próprio autor'' ou ``elaboração própria'' como fonte.

Para que a legenda fique sempre acima do conteúdo, como a norma exige, use os ambientes
\texttt{figuraabnt} e \texttt{tabelaabnt} deste modelo em vez dos ambientes \texttt{figure} e
\texttt{table} do \LaTeX. Neles, o rótulo e a legenda são argumentos do próprio ambiente, e não
comandos escritos no meio do corpo, de modo que não há ordem errada possível de digitar:

\begin{verbatim}
\begin{figuraabnt}[!htbp]{fig:rotulo}{Legenda da figura}[Legenda curta]
    \includegraphics[width=0.9\textwidth]{arquivo}
    \fonte{o próprio autor}
\end{figuraabnt}
\end{verbatim}

O primeiro argumento, entre colchetes, é a posição do float e pode ser omitido (o padrão é
\texttt{[!htbp]}). Seguem o rótulo e a legenda, ambos obrigatórios. O próximo argumento, também entre
colchetes e opcional, é uma legenda curta, usada apenas na Lista de Ilustrações quando a legenda
principal for longa demais para caber ali. O ambiente \texttt{tabelaabnt} tem exatamente a mesma
assinatura. Os ambientes \texttt{figure} e \texttt{table} continuam disponíveis, mas neles a posição
da legenda volta a depender da ordem em que você escreve o código, sem nenhum aviso caso saia fora
da norma.

Há ainda um último argumento opcional, entre \texttt{<} e \texttt{>}: a largura declarada da
ilustração. A norma exige que ``tipo, número de ordem, título, fonte, legenda e notas devem
acompanhar as margens da ilustração'' (§5.8) — ou seja, quando o gráfico é mais estreito que o bloco
de texto, a legenda e a fonte devem se alinhar à largura do gráfico, não à largura da página. Sem
esse argumento, a legenda ocupa o bloco de texto inteiro e a fonte se alinha à largura medida do
último gráfico incluído, como sempre foi. Repita nele o mesmo valor passado a
\texttt{\textbackslash includegraphics}:

\begin{verbatim}
\begin{figuraabnt}{fig:rotulo}{Legenda}<0.9\textwidth>
    \includegraphics[width=0.9\textwidth]{arquivo}
    \fonte{o próprio autor}
\end{figuraabnt}
\end{verbatim}

Veja aqui uma referência à \reffig{fig:exemplo} e à \reftable{tab:exemplo}.

\begin{figuraabnt}{fig:exemplo}{Figura de exemplo}<0.9\textwidth>
    \caption*{Uma figura de exemplo, elaborada pelo próprio autor. Note o alinhamento da ``fonte'', abaixo da figura. Esta legenda adicional é opcional, usada quando é necessário fornecer uma descrição mais detalhada da figura, complementando a legenda principal. Note também que a legenda e a fonte acompanham a largura declarada da figura.}
    \includegraphics[width=0.9\textwidth]{example-image}
    \fonte{o próprio autor}
\end{figuraabnt}

\begin{tabelaabnt}{tab:exemplo}{Tabela de exemplo}
    \caption*{A legenda da tabela deve conter uma descrição clara e concisa do conteúdo apresentado.}
    \begin{tabular}{@{}lllll@{}}
        \toprule
        Coluna 1 & Coluna 2 & Coluna 3 & Coluna 4 & Coluna 5 \\ \midrule
        1        & 2        & 2        & 1        & 3        \\
        1        & 2        & 1        & 4        & 3        \\
        5        & 3        & 4        & 3        & 5        \\ \bottomrule
    \end{tabular}
    \fonte{o próprio autor}
\end{tabelaabnt}

Usando o pacote \texttt{subcaption}, é possível incluir múltiplas figuras ou tabelas dentro de um
mesmo ambiente \texttt{figuraabnt} ou \texttt{tabelaabnt}, cada uma com sua própria legenda, além da
legenda principal --- esta última continua sendo o argumento do ambiente. Veja a \reffig{fig:exemplo-multiplo} como exemplo.

\begin{figuraabnt}{fig:exemplo-multiplo}{Exemplo de figura múltipla}[Figura múltipla]
    \begin{subfigure}{0.45\textwidth}
        \centering
        \includegraphics[width=\textwidth]{example-image-a}
        \caption{Subfigura A}
        \label{fig:subfig-a}
    \end{subfigure}
    \begin{subfigure}{0.45\textwidth}
        \centering
        \includegraphics[width=\textwidth]{example-image-b}
        \caption{Subfigura B}
        \caption*{Cada subfigura pode ter sua própria legenda}
        \label{fig:subfig-b}
    \end{subfigure}
    \fonte{(a): \cite{ibge2020}; (b): \cite{almeida2022}}
\end{figuraabnt}

Lembre-se que você pode sempre consultar a documentação do pacote \texttt{subcaption} para mais
detalhes sobre como utilizá-lo de forma eficaz. Ele pode ser encontrado no diretório oficial do
\LaTeX, na \gls{ctan}, no endereço eletrônico (\url{https://ctan.org/pkg/subcaption}).

\subsection{Uma limitação do \texttt{tabular} neste modelo}\label{ssc:verb-tabular}

\textbf{\texttt{\textbackslash verb} não funciona dentro de um \texttt{tabular} neste modelo.}
Escrevê-lo numa célula interrompe a compilação com \texttt{\textbackslash verb illegal in
argument}, e a mensagem aponta para o \texttt{\textbackslash end\{tabular\}} --- longe da célula
que a causou.

Não é defeito do \LaTeX{}, e sim consequência de o modelo redefinir o \texttt{tabular} para medir a
largura da tabela: o corpo passa a ser lido como argumento, e \texttt{\textbackslash verb} é
ilegal em argumento, por construção. Nada disso é visível de fora, e a limitação vale para toda
tabela deste modelo.

Para escrever código numa célula, use o \texttt{\textbackslash texttt} com o contrabarra
explícito:

\begin{Verbatim}[fontsize=\small]
\texttt{\textbackslash comando}     % produz: \comando
\end{Verbatim}

\noindent Fora de tabelas, o \texttt{\textbackslash verb} continua disponível como sempre.

\subsection{Outros tipos de ilustração}\label{sec:tipos-ilustracao}

A norma não fala apenas em ``figura''. O §5.8 exige que toda ilustração venha precedida de sua
palavra designativa específica — desenho, esquema, fluxograma, fotografia, gráfico, mapa,
organograma, planta, quadro, retrato, figura, imagem, entre outros —, cada uma com sua própria
sequência numérica e, quando conveniente, sua própria lista (§4.2.1.9). Além de \texttt{figuraabnt}
e \texttt{tabelaabnt}, este modelo já fornece \texttt{quadroabnt} e \texttt{graficoabnt}, com
exatamente a mesma assinatura — incluindo a largura declarada opcional \texttt{<largura>}. Veja o
\refgrafico{gra:exemplo} como exemplo:

\begin{graficoabnt}{gra:exemplo}{Gráfico de exemplo}<0.7\textwidth>
    \includegraphics[width=0.7\textwidth]{example-image}
    \fonte{o próprio autor}
\end{graficoabnt}

Um relatório que precise de um tipo não previsto aqui (Fluxograma, Organograma, ...) não precisa
esperar uma atualização do pacote. Basta declarar o tipo no preâmbulo do próprio documento:

\begin{verbatim}
\novotipoilustracao{fluxograma}{Fluxograma}{fluxogramaabnt}
\end{verbatim}

Os três argumentos são o nome interno do contador (só ASCII), a palavra designativa como deve
aparecer impressa, e o nome do ambiente que você vai usar no texto. A partir daí,
\texttt{fluxogramaabnt} funciona exatamente como \texttt{quadroabnt} ou \texttt{graficoabnt} --- com
contador, referências cruzadas (\verb|\reffluxograma{}|) e lista (\verb|\listoffluxogramas|)
próprios. É a mesma máquina que gera os tipos que já vêm no pacote.

Cada tipo tem o par completo de comandos de referência, como Figura e Tabela: \texttt{\textbackslash refquadro}
e \texttt{\textbackslash refgrafico} citam o número, e \texttt{\textbackslash refquadrocomp} e
\texttt{\textbackslash refgraficocomp} citam o número junto do título. Um tipo criado por
\texttt{\textbackslash novotipoilustracao} ganha os dois na mesma hora.

\subsubsection{Quadro não é Tabela}\label{sec:quadro-vs-tabela}

Este é o erro mais comum: chamar de ``Tabela'' qualquer coisa organizada em linhas e colunas. A
norma traça a fronteira pelo tipo de dado central, não pelo layout. Uma tabela é a forma não
discursiva de apresentar informações das quais o \emph{dado numérico} se destaca como informação
central (§5.9, vinculado às normas de apresentação tabular do IBGE); um quadro é uma ilustração
comum (§5.8) cujo conteúdo central é \emph{textual}, ainda que organizado em linhas e colunas. Veja
o \refquadro{qua:exemplo} e a \reftable{tab:exemplo-quadro} a seguir: é a mesma avaliação, uma vez
qualitativa (Quadro) e outra numérica (Tabela).

\begin{quadroabnt}{qua:exemplo}{Quadro de exemplo}
    \caption*{Avaliação qualitativa: o dado central de cada célula é uma palavra, não um número — por isso é um Quadro, não uma Tabela.}
    \begin{tabular}{@{}ll@{}}
        \toprule
        Critério    & Avaliação \\ \midrule
        Usabilidade & Boa       \\
        Desempenho  & Regular   \\
        Documentação & Boa      \\
        Suporte     & Ruim      \\ \bottomrule
    \end{tabular}
    \fonte{o próprio autor}
\end{quadroabnt}

\begin{tabelaabnt}{tab:exemplo-quadro}{Tabela de exemplo (mesma avaliação, em números)}
    \caption*{A mesma avaliação do \refquadro{qua:exemplo}, agora com o dado numérico como informação central — por isso é uma Tabela.}
    \begin{tabular}{@{}ll@{}}
        \toprule
        Critério    & Nota (0--10) \\ \midrule
        Usabilidade & 8            \\
        Desempenho  & 6            \\
        Documentação & 8           \\
        Suporte     & 3            \\ \bottomrule
    \end{tabular}
    \fonte{o próprio autor}
\end{tabelaabnt}

\section{Equações}\label{sec:equacoes}

Equações também são elementos comuns em trabalhos acadêmicos, especialmente em áreas como
matemática, física e engenharia. Equações podem ser apresentadas de duas formas, \emph{inline} ou
em blocos separados. Equações \emph{inline} são inseridas diretamente no texto, enquanto equações em
blocos separados são destacadas do texto, geralmente centralizadas e numeradas para referência. Veja
um exemplo de equação \emph{inline}, usando a equação considerada por muitos, como a mais bela de
todas, a \emph{Identidade de Euler}: \(e^{i\pi} + 1 = 0\). O \LaTeX\ provê uma maneira simples de
formatar equações matemáticas, com uma variedade de símbolos e operadores disponíveis. Para equações
em blocos separados, utilize o ambiente \texttt{equation}, como mostrado a seguir:

\begin{equation}
    a^2 + b^2 = c^2
    \label{eq:pitagoras}
\end{equation}

Equações em blocos também podem ser referenciadas no texto. O modelo provê os comandos
\verb|\refeq{}| e \verb|\refeqcomp{}| para referenciar equações com formatação consistente com o
estilo do documento. Por exemplo, veja a referência à \refeq{eq:pitagoras}. E a referência
completa à mesma: \refeqcomp{eq:pitagoras}[Equação do Teorema de Pitágoras]. Note que o nome
customizado vai entre colchetes, por ser um argumento opcional; sem ele, basta
\verb|\refeqcomp{eq:pitagoras}|.

\section{Nomes que você pode trocar}\label{sec:nomes-personalizaveis}

Duas palavras que o modelo imprime sozinho vêm de comandos que você pode redefinir no preâmbulo:

\begin{itemize}
    \item \texttt{\textbackslash figuresourcename} --- a palavra que precede a fonte de uma
          ilustração. Padrão em português: \textit{Fonte}.
    \item \texttt{\textbackslash eqname} --- a palavra designativa das equações, usada por
          \texttt{\textbackslash refeq} e \texttt{\textbackslash refeqcomp}. Padrão:
          \textit{Equação}.
\end{itemize}

As duas seguem o idioma declarado no \texttt{babel} --- em espanhol saem \textit{Fuente} e
\textit{Ecuación}, e assim por diante. Redefina apenas se precisar de algo que o idioma não dá:

\begin{Verbatim}[fontsize=\small]
\renewcommand{\figuresourcename}{Elaborado a partir de}
\end{Verbatim}

A redefinição vale do ponto em que aparece em diante, então o lugar dela é o preâmbulo.
